% Generated draft: requires mathematical and source review.
\documentclass[11pt]{article}
\usepackage[margin=25mm]{geometry}
\usepackage{fontspec}
\setmainfont{Latin Modern Roman}
\usepackage{amsmath,amssymb,mathtools,unicode-math}
\setmathfont{Latin Modern Math}
\usepackage{enumitem,needspace,xurl,hyperref}
\setlength{\parindent}{0pt}
\setlength{\parskip}{5pt}
\setlength{\emergencystretch}{3em}
\newcommand{\N}{\mathbb{N}}
\newcommand{\Z}{\mathbb{Z}}
\newcommand{\Q}{\mathbb{Q}}
\newcommand{\R}{\mathbb{R}}
\newcommand{\C}{\mathbb{C}}
\newcommand{\F}{\mathbb{F}}
\newcommand{\A}{\mathbb{A}}
\newcommand{\PP}{\mathbb{P}}
\newcommand{\E}{\mathbb{E}}
\newcommand{\eps}{\varepsilon}
\newcommand{\dd}{\,\mathrm{d}}
\begin{document}
\section*{UnsolvedMath ID: 30000672}
\textbf{Source problem code:} catalogue entry
\subsection{Definitions and mathematical statement}
\textbf{Relative Galois-reconstruction question, curated form.} Let \(K,K'\) be perfect fields. Assume \(K\) has a finite extension of degree \(d>2\) prime to \(\operatorname{char}K\). Let \(F/K\) be a one-variable function field, and let \(F'/K'\) be any function field. Is the canonical map\[\operatorname{Isom}(F/K,F'/K')\longrightarrow\operatorname{Isom}(G_{F/K},G_{F'/K'})/\operatorname{Inn}(G_{F'/K'})\] bijective? A perfect field has all algebraic extensions separable, equivalently surjective Frobenius in positive characteristic. In characteristic \(p>0\), the degree restriction means \(p\nmid d\); in characteristic zero, the intended prime-to-characteristic condition conventionally imposes no additional divisibility restriction. The source does not state that convention explicitly. The finite extension is not required to be Galois by the supplied wording.\par One-variable means transcendence degree one for the intended finitely generated extension. The source does not fix regularity or constant-field conventions in \emph{function field}. The curated version does not restrict the transcendence degree of \(F'/K'\) to one, so that broader quantifier is retained. A conventional source-side isomorphism consists of an isomorphism \(F\to F'\) carrying \(K\) onto \(K'\), with its restriction furnishing the base-field isomorphism; whether bases are instead preidentified is not specified.\par The original calls \(G_{F/K}\) a corresponding relative absolute Galois group but does not define it. A geometric Galois kernel, an absolute Galois group equipped with a map to the constant-field Galois group, and another relative automorphism construction are different data. They must not be interchanged from this symbol alone. Its topology, required compatibility with constant-field structure, and the category of group isomorphisms are missing. The target quotient is by postcomposition with inner automorphisms of the target group. A standard absolute-Galois construction uses extensions of field isomorphisms to separable closures and removes the choice by this quotient, but that description does not identify the missing relative construction. Bijectivity asks both recovery from every permitted outer isomorphism and uniqueness of the field isomorphism.\par The original prose instead restricts both function fields to one variable and omits the degree-\(d\) hypothesis on \(K\). The later curated version adds that hypothesis and broadens \(F'/K'\), while its verification says source checking preserved the claim. These materially different source layers are retained explicitly; no equivalence is inferred without the actual report definitions.\par The supplied triage checked 2026-08-21 labels the extracted statement open and reports no located complete resolution, citing only the report. It supplies no precise later theorem or reconstruction category and requires independent verification.
\subsection{Short English statement}
Ask whether relative field isomorphisms are exactly the permitted outer isomorphisms of the corresponding relative Galois objects. The Galois construction and compatibility category are undefined. The curated extra base-field hypothesis and unrestricted target transcendence degree differ from the original prose and remain explicit.
\subsection{Sources}
\begin{itemize}\item Local UnsolvedMath registry, permanent ID 30000672, code catalogue entry; \emph{Relative Galois Reconstruction of One-Variable Function Fields}, Oberwolfach Reports Open Problems.\item \emph{Model Theory and Groups (2007)}, Oberwolfach Reports, DOI 10.4171/owr/2007/02, pp. 83--138. \url{https://doi.org/10.4171/owr/2007/02}. Proposed year 2007; no proposer named.\item The same full report is the sole supplied literature source. No exact definition of the relative Galois group or separate authoritative reconstruction theorem is supplied.\item Literature triage checked 2026-08-21, status \texttt{open}; record updated 2026-09-16. The \texttt{corrected\_verified} field reports source checking, not independently verified here.\end{itemize}
\end{document}
